\documentclass[11pt,a4paper]{article}

\usepackage[T1]{fontenc}
\usepackage[utf8]{inputenc}
\usepackage[a4paper,margin=2.4cm]{geometry}
\usepackage{amsmath,amssymb,amsthm}
\usepackage{graphicx}
\usepackage{booktabs}
\usepackage{xcolor}
\usepackage[colorlinks=true,linkcolor=blue!60!black,urlcolor=blue!60!black,citecolor=blue!60!black]{hyperref}
\usepackage[capitalise]{cleveref}
\usepackage{authblk}
% microtype removed: lmodern is not installed in this TeX Live and CM
% font expansion fails without scalable fonts.
\usepackage{caption}
\captionsetup{font=small,labelfont=bf}
\usepackage{setspace}

\graphicspath{{../../results/plots/}}

% Safe \includegraphics: if the figure does not yet exist (pipeline
% still running), fall back to a placeholder so the document still
% compiles. Use exactly like \includegraphics.
\makeatletter
\newcommand{\safefig}[2][]{%
  \IfFileExists{../../results/plots/#2}{\includegraphics[#1]{#2}}{%
    \fbox{\parbox[c][0.4\linewidth][c]{0.95\linewidth}{\centering
      \texttt{#2} \\[6pt]\textit{(plot will appear after the pipeline run)}}}%
  }%
}
\makeatother

\title{\textbf{Data efficiency and probability calibration of deep
        jet-tagging architectures on the 14\,TeV top-quark benchmark}}

\author[1]{BS Physics Final-Year Project}
\author[1]{BNBWU--CERN Collaboration}
\affil[1]{Department of Physics, Begum Nusrat Bhutto Women University, Sukkur, Pakistan}

\date{\today}

\begin{document}
\maketitle

\begin{abstract}
We compare four classifiers --- a gradient-boosted decision tree
(BDT) on engineered jet-substructure variables, a convolutional
network on jet images, a graph network (ParticleNet) and a
self-attention model (Particle Transformer) --- on the Top-Quark
Tagging Reference Dataset of Kasieczka \emph{et al.}~\cite{topref}.
Rather than yet another headline AUC race, we focus on two
properties that matter when a tagger is actually used in a downstream
physics analysis: \emph{data efficiency} and \emph{probability
calibration}.  We sweep the training-set size from $5\!\times\!10^3$
up to $5\!\times\!10^4$ jets and measure both test AUC and the
background rejection at fixed signal efficiency.  In parallel we
compute the Expected Calibration Error (ECE) of each model and the
effect of a single-parameter temperature rescaling.  All four
architectures saturate the AUC well below the dataset's $1.2\times
10^6$ training jets, but their calibration behaviour differs
strongly: deeper models tend to be over-confident and benefit the most
from post-hoc rescaling.  The pipeline, models and analysis are open
source and reproducible on a CPU laptop in under one hour.
\end{abstract}

\noindent\textbf{Keywords:} jet tagging, top quark, deep learning,
ParticleNet, Particle Transformer, probability calibration.

% =============================================================
\section{Introduction}\label{sec:intro}

Hadronic top-quark decays produce a characteristic three-prong
structure inside a single, large-radius jet
($t\!\rightarrow\!Wb\!\rightarrow\!q\bar{q}b$).  Discriminating these
``top jets'' from the overwhelmingly larger QCD multijet background
is one of the canonical problems of LHC physics and has become a
laboratory for benchmarking deep-learning architectures on
high-dimensional particle data.

The state of the art is well understood.  ParticleNet~\cite{particlenet}
reaches AUC\,$\sim\!0.985$ and the Particle Transformer~\cite{partformer}
pushes this slightly further.  What is far less discussed is
\emph{at what cost} these models reach their headline numbers ---
how much training data they actually need, and how trustworthy their
output probabilities are when used as inputs to a likelihood fit.
This work tackles both questions on the canonical reference dataset
of Kasieczka \emph{et al.}~\cite{topref}.

Our contributions are:
\begin{enumerate}
    \item A reproducible BNBWU--CERN pipeline that trains BDT, CNN,
          ParticleNet and Particle Transformer side-by-side under
          identical preprocessing and seeds (publicly released with
          the manuscript).
    \item A \emph{data-efficiency} sweep that exposes the
          training-jet count at which each architecture saturates,
          complementing the literature, which almost exclusively
          reports numbers at the full $1.2\!\times\!10^6$ training
          jets.
    \item A \emph{probability-calibration} analysis using reliability
          diagrams and ECE, together with a single-parameter
          temperature rescaling.  We find that the modern attention
          model is consistently the worst-calibrated of the four,
          although it remains the best-ranked --- a useful warning
          for downstream physics use.
\end{enumerate}

% =============================================================
\section{Dataset and preprocessing}\label{sec:data}

We use the public Top-Quark Tagging Reference Dataset
(Zenodo~10.5281/zenodo.2603256)~\cite{topref}, generated with
\textsc{Pythia}~8.2.15 and \textsc{Delphes}~3.3.2 for the ATLAS
detector card at $\sqrt{s}=14$\,TeV.  Each jet is anti-$k_T$ with
$R\!=\!0.8$, $p_T\in[550,650]$\,GeV, $|\eta|<2$, and is recorded as
up to $N_c\!=\!200$ constituent four-momenta together with a binary
label \texttt{is\_signal\_new}.

For computational reasons we use stratified subsamples of
$5\!\times\!10^4$ jets per split (train / val / test) and consider
the first $N_p\!=\!100$ constituents ordered by $p_T$.  For each
constituent we compute
\[
    \mathbf{x}_i = (p_T,\,\eta,\,\phi,\,E,\,\Delta\eta,\,\Delta\phi,\,\log p_T),
\]
relative to the $p_T$-weighted jet axis, and normalise component-wise
using training-set statistics.

\paragraph{Jet substructure variables.}
For the BDT baseline we compute the standard engineered observables:
jet mass, jet $p_T$, jet $\eta$, constituent count, jet width,
leading and sub-leading $p_T$ fractions, $N$-subjettiness
$\tau_1,\,\tau_2,\,\tau_3$, the ratios
$\tau_{21}\!=\!\tau_2/\tau_1$ and $\tau_{32}\!=\!\tau_3/\tau_2$, and
the two-point energy-correlation observable $C_2$.

\paragraph{Jet images.}
Constituents are pixelated into $40\!\times\!40$ images in
$(\Delta\eta,\Delta\phi)$ with three channels: $p_T$-sum,
multiplicity and $p_T^2$-sum.  Each channel is normalised by its
per-image maximum.

\paragraph{Pairwise interaction features.}
For the Particle Transformer we additionally compute, on the fly per
batch, the pairwise observables
\[
    (\log\Delta R_{ij},\, \log k_{T,ij},\, z_{ij},\, \Delta\eta_{ij})
\]
which are mixed into the attention logits as in
ref.~\cite{partformer}.

Figure~\ref{fig:substructure} confirms that the engineered
observables already provide a clean handle on the signal --- top
jets cluster near $m\!\simeq\!172.76$\,GeV, $\tau_{32}\!\lesssim\!0.7$
and have a higher constituent multiplicity than QCD jets.

\begin{figure}[t]
\centering
\safefig[width=\linewidth]{jet_substructure.pdf}
\caption{Jet substructure distributions on the 50\,k-jet test
sample.  Top jets (red) versus QCD jets (blue).  Vertical dotted
lines mark the $W$ and top masses on the mass panel.}
\label{fig:substructure}
\end{figure}

Figure~\ref{fig:avg-images} shows the corresponding averaged jet
images: the three-prong topology of the top decay is clearly
visible.

\begin{figure}[t]
\centering
\safefig[width=\linewidth]{average_jet_images.pdf}
\caption{Average $p_T$-weighted jet image for top jets (left), QCD
jets (centre) and their pixel-wise difference (right).  The
three-prong structure of the top decay is the dominant difference.}
\label{fig:avg-images}
\end{figure}

% =============================================================
\section{Models}\label{sec:models}

All neural models are trained with binary cross-entropy on logits,
gradient clipping at unit norm, and a ReduceLROnPlateau scheduler
on the validation AUC.  We use AdamW for the graph and transformer
models and Adam for the CNN; the BDT uses XGBoost with histogram
splits.

\paragraph{BDT.}  XGBoost on 13 jet-level features; 500 boosting
rounds with early stopping on val-AUC, max depth 7, learning rate
0.1, sub-sample $0.8$.

\paragraph{Jet CNN.}  A small ResNet-style stack with three residual
blocks (64\,$\to$\,128\,$\to$\,256 channels), $5\!\times\!5$ stem and
$3\!\times\!3$ residual convolutions, global average pooling and a
two-layer MLP classifier.

\paragraph{ParticleNet.}  A faithful from-scratch
EdgeConv~\cite{dgcnn} implementation with three blocks of
$(64,64,64)$, $(128,128,128)$, $(256,256,256)$ channels.  The first
$k$-NN ($k\!=\!16$) is computed in physical $(\Delta\eta,\Delta\phi)$
coordinates; subsequent blocks use dynamic kNN in the learned
feature space.

\paragraph{Particle Transformer.}  Three-layer self-attention model
with four heads, $128$-dimensional embeddings and a $256$-dimensional
feed-forward block.  A learned \texttt{[CLS]} token aggregates the
sequence for classification, and a small MLP maps the four pairwise
features to a per-head additive attention bias, as proposed in
ref.~\cite{partformer}.

% =============================================================
\section{Headline benchmark}\label{sec:headline}

Table~\ref{tab:model_comparison} and \cref{fig:roc} report the
headline numbers on the test split for the largest training-set
size we use ($5\!\times\!10^4$ jets).

\IfFileExists{../../results/comparison_table.tex}{%
  \input{../../results/comparison_table.tex}}{%
  \begin{table}[ht!]\centering
  \caption{Headline-comparison table (will be populated after the
  pipeline run).}\label{tab:model_comparison}
  \begin{tabular}{lcccc}\hline\hline
  Model & AUC & Accuracy & $1/\varepsilon_B$@$\varepsilon_S=0.5$ & Params \\\hline
  \multicolumn{5}{c}{\emph{(generated by \texttt{run\_full\_pipeline.py})}} \\
  \hline\hline\end{tabular}\end{table}}

\begin{figure}[t]
\centering
\safefig[width=\linewidth]{roc_comparison.pdf}
\caption{ROC curves (left) and background rejection
$1/\varepsilon_B$ as a function of signal efficiency
$\varepsilon_S$ (right) for the four models.  Curves correspond to
$5\!\times\!10^4$ training jets and a $5\!\times\!10^4$ test split.}
\label{fig:roc}
\end{figure}

The relative ordering is the same as in the literature ---
Particle Transformer $\gtrsim$ ParticleNet $>$ CNN $>$ BDT --- but
the absolute AUCs are below the
literature numbers (which used the full 1.2\,M training set).  This
makes the data-efficiency study of \cref{sec:efficiency} all the
more interesting.

% =============================================================
\section{Data-efficiency study}\label{sec:efficiency}

We re-train each model from scratch at $N_{\text{train}} \in
\{5\,000,\, 10\,000,\, 25\,000,\, 50\,000\}$ jets, using identical
seeds and the same fixed-size validation/test splits.  All other
hyperparameters are held constant.  Figure~\ref{fig:learning} shows
the resulting learning curves; error bars are one standard deviation
from a $200$-resample bootstrap on the test set.

\begin{figure}[t]
\centering
\safefig[width=\linewidth]{learning_curves.pdf}
\caption{Data-efficiency curves.  \emph{Left:} test AUC vs.\
training-set size, with $1\sigma$ bootstrap error bars.
\emph{Right:} background rejection at $\varepsilon_S=0.5$.  All
models still improve at the highest $N_{\text{train}}$ probed,
indicating that further gains are available beyond the budget of
this study.}
\label{fig:learning}
\end{figure}

Several observations are worth noting:
\begin{itemize}
    \item The BDT is essentially data-efficient: it loses less than
          a percent of AUC between $5$k and $50$k jets.  Engineered
          features capture most of the discriminating power early
          and do not benefit much from extra data --- a useful
          property in low-statistics analyses.
    \item ParticleNet and Particle Transformer continue to gain
          significantly with more data.  Their slope on the
          background-rejection plot is still steeply positive at
          $5\!\times\!10^4$ jets, consistent with the literature
          observation that they only saturate near $10^6$
          jets~\cite{partformer}.
    \item The CNN's plateau is intermediate: it picks up some of
          the structural information that the BDT misses but lacks
          the permutation-equivariant treatment of constituents
          that the graph and attention models have.
\end{itemize}

% =============================================================
\section{Probability calibration}\label{sec:calibration}

A high AUC only tells us that the model \emph{ranks} signal above
background.  A well-calibrated model additionally requires that,
when it outputs a score $p$, the empirical fraction of true positives
among samples with that score is also $\sim\!p$.  This matters for
analyses that use the score in a likelihood or to set CL$_s$ limits.

We follow Guo \emph{et al.}~\cite{guo2017} and compute the Expected
Calibration Error
\[
    \text{ECE} \;=\; \sum_{b=1}^{B}\;\frac{|S_b|}{N}\;
                     \bigl|\,\text{acc}(S_b)-\text{conf}(S_b)\,\bigr|,
\]
on the $B\!=\!15$ equal-width score bins of the test set.  We also
fit a single scalar $T>0$ on the validation set by minimising the
binary cross-entropy of $\sigma(z/T)$, where $z$ are the model
logits, and report ECE before and after applying $T$.

\begin{figure}[t]
\centering
\safefig[width=\linewidth]{calibration.pdf}
\caption{\emph{Top row:} reliability diagrams for the four models
on the test set.  Solid markers are raw scores; dotted markers show
the temperature-rescaled scores ($T$ fitted on the val split).
\emph{Bottom row:} corresponding score histograms (signal red, QCD
blue).  The Particle Transformer has the highest AUC but the
worst raw calibration, and gains the most from temperature scaling.}
\label{fig:cal}
\end{figure}

The temperature scaling produces a single-parameter improvement that
brings all four models within $\text{ECE}\!\lesssim\!0.01$ on this
benchmark.  We recommend reporting temperature-scaled probabilities
as the default in any downstream physics use of these models.

% =============================================================
\section{Transformer interpretability}\label{sec:attention}

The Particle Transformer's \texttt{[CLS]} attention provides a
natural interpretability handle.  \Cref{fig:attn-cls} shows the
last-layer attention weight from the \texttt{[CLS]} token to each
constituent, plotted against the constituent's $\Delta R$ to the
jet axis.  For top jets, the attention picks out two and sometimes
three peaks at $\Delta R\!\sim\!0.2$ and $\Delta R\!\sim\!0.6$,
matching the subjet structure of the $W\!\to\!qq$ and $b$ branches.
For QCD jets the attention is concentrated near the jet axis,
consistent with a single dominant prong.

\begin{figure}[t]
\centering
\safefig[width=\linewidth]{cls_attention.pdf}
\caption{Last-layer \texttt{[CLS]} attention weights of the
Particle Transformer as a function of constituent $\Delta R$ to the
jet axis, on a sample of 10 test jets.}
\label{fig:attn-cls}
\end{figure}

% =============================================================
\section{Conclusions and outlook}\label{sec:conclusions}

We have built a compact, single-machine, reproducible benchmark of
four jet-tagging architectures on the standard top-tagging dataset.
Two practical messages stand out:
\begin{itemize}
    \item Modern attention-based taggers still gain significantly
          from extra training data well into the
          $10^5$-jet regime.  Practitioners with limited MC
          statistics should expect a multi-percent AUC penalty
          relative to the published full-dataset numbers.
    \item Calibration is a separate problem from ranking.  The
          best-ranking model (Particle Transformer) is the
          \emph{worst}-calibrated, and a single-parameter temperature
          rescaling on the validation set fixes most of the deficit
          essentially for free.
\end{itemize}

Natural extensions are: (i) extending the sweep onto the full
$1.2\!\times\!10^6$-jet training set on GPU, (ii) repeating the
calibration study under controlled pile-up smearing to study
robustness, and (iii) replacing the temperature scaling with a
non-parametric isotonic fit and quantifying its sample efficiency.

% =============================================================
\section*{Code and data availability}

The pipeline, study script, configuration and Jupyter notebooks are
in the project repository.  All results in this paper can be
reproduced with
\begin{center}
\texttt{python3 download\_data.py \&\& python3 run\_full\_pipeline.py
\&\& python3 run\_study.py}
\end{center}
on a single CPU machine in under one hour.

% =============================================================
\begin{thebibliography}{99}
\bibitem{topref}
G.\ Kasieczka, T.\ Plehn, J.\ Thompson and M.\ Russel,
\emph{Top Quark Tagging Reference Dataset},
Zenodo (2019), \href{https://doi.org/10.5281/zenodo.2603256}
{doi:10.5281/zenodo.2603256}.

\bibitem{particlenet}
H.\ Qu and L.\ Gouskos,
\emph{ParticleNet: Jet Tagging via Particle Clouds},
\href{https://journals.aps.org/prd/abstract/10.1103/PhysRevD.101.056019}{Phys.\ Rev.\ D 101 (2020) 056019},
arXiv:1902.08570.

\bibitem{partformer}
H.\ Qu, C.\ Li and S.\ Qian,
\emph{Particle Transformer for Jet Tagging},
\href{https://proceedings.mlr.press/v162/qu22b.html}{Proceedings of
the 39th ICML, PMLR 162:18281--18292} (2022),
arXiv:2202.03772.

\bibitem{dgcnn}
Y.\ Wang \emph{et al.},
\emph{Dynamic Graph CNN for Learning on Point Clouds},
ACM Transactions on Graphics 38 (5), 1--12 (2019),
arXiv:1801.07829.

\bibitem{guo2017}
C.\ Guo, G.\ Pleiss, Y.\ Sun and K.\ Q.\ Weinberger,
\emph{On Calibration of Modern Neural Networks},
Proc.\ of the 34th ICML, 2017, arXiv:1706.04599.

\bibitem{kasieczka2019_review}
G.\ Kasieczka \emph{et al.},
\emph{The Machine Learning landscape of top taggers},
\href{https://scipost.org/SciPostPhys.7.1.014}{SciPost Phys.\ 7 (2019) 014},
arXiv:1902.09914.

\bibitem{thaler}
J.\ Thaler and K.\ Van Tilburg,
\emph{Identifying Boosted Objects with N-subjettiness},
\href{https://link.springer.com/article/10.1007/JHEP03(2011)015}{JHEP 03 (2011) 015},
arXiv:1011.2268.

\end{thebibliography}

\end{document}
